\subsection*{B.4 Entrevista N.°~4 --- Julián Fabbi}
\addcontentsline{toc}{subsection}{B.4 Entrevista N.° 4 --- Julián Fabbi}

\noindent\textbf{Entrevistado:} Julián Fabbi\\
\noindent\textbf{Perfil:} Dueño de emprendimiento de venta de cómics, historietas y mangas. Vende por Tienda Nube y Mercado Libre. No tiene tienda física. Rubro regulado por la Ley 25.542 de protección de la actividad librera (precios fijados por editorial). Desde octubre de 2025.\\
\noindent\textbf{Entrevistador:} Tomás Gonçalves (moderador, sin observador en esta sesión)\\
\noindent\textbf{Fecha:} Jueves 24 de julio de 2026\\
\noindent\textbf{Modalidad:} Videollamada (semiestructurada) --- Duración: \textasciitilde28 minutos\\
\noindent\textbf{Plataforma del negocio:} Tienda Nube (plan básico, \$27.000/mes) + Mercado Libre

\begin{framed}
\small\noindent\textit{Nota metodológica: Esta transcripción fue limpiada a partir de una grabación de \textasciitilde28 minutos vía Microsoft Teams. Se eliminó charla trivial y se corrigió la sintaxis para mejorar la fluidez de lectura. Se filtró una sección extensa sobre la Ley 25.542 de protección librera, conservando únicamente lo que impacta directamente en la gestión de precios del entrevistado. No se agregó información que el entrevistado no haya dicho explícitamente.}
\end{framed}

\subsubsection*{Bloque 0 · Introducción y rapport}

\noindent\textbf{Tomás Gonçalves:} Nosotros estamos haciendo nuestro proyecto final: una herramienta que ayuda a dueños de tienda online a tomar mejores decisiones usando sus propios datos. Nosotros le decimos recomendaciones accionables. La idea es que agarra directamente tus datos y en base a eso te dice, por ejemplo: el martes que viene podés poner esta promoción y tenés asegurada esta cantidad de ventas.

\noindent\textbf{Tomás Gonçalves:} Entonces, para empezar: ¿qué vendés y por dónde vendés?

\noindent\textbf{Julián Fabbi:} Vendo cómics e historietas: las que vienen de Estados Unidos (las más conocidas, Marvel y DC), acá en Argentina también hay una gran industria historietista en la que me estoy empezando a meter hace poquito, y lo principal que vendo son los mangas, que son los cómics japoneses. Vendo a través de dos canales virtuales: uno es Tienda Nube y el otro es Mercado Libre. De momento no tengo tienda física.

\noindent\textbf{Tomás Gonçalves:} ¿Cuándo arrancaste más o menos?

\noindent\textbf{Julián Fabbi:} En octubre del año pasado, 2025. El 10 o 15 de octubre arranqué.

\subsubsection*{Bloque 1 · Contexto del negocio y herramientas actuales}

\noindent\textit{--- Registro de ventas y uso de plataformas ---}

\noindent\textbf{Tomás Gonçalves:} Para registrar las ventas, ¿usás Tienda Nube, Excel, alguna app?

\noindent\textbf{Julián Fabbi:} Tienda Nube registra todo: tiene el registro de stock, o sea que cada vez que compran, descuenta directamente. Lo que tengo que hacer yo es subirlo. Tienda Nube también registra los ingresos y tiene un pequeño apartado estadístico, que requiere un nivel más alto de pago y no le presto tanta atención. Aparte de eso, llevo los registros en Excel porque hay algunos valores que Tienda Nube no contempla.

\noindent\textbf{Tomás Gonçalves:} ¿Le das mucha pelota? ¿Lo mirás con frecuencia o cada fin de mes?

\noindent\textbf{Julián Fabbi:} Cada fin de mes veo todo y lo voy guardando. Lo llevo en un registro y estoy atento a cómo vienen ingresos y egresos para tantear las compras que vaya a hacer semanalmente. Uso Excel porque en Excel consolido los ingresos de ambas plataformas.

\noindent\textit{--- Métricas de Tienda Nube y datos faltantes ---}

\noindent\textbf{Tomás Gonçalves:} Tienda Nube te genera cosas. ¿Hay algo que te gustaría saber y que no te está mostrando?

\noindent\textbf{Julián Fabbi:} Me muestra la gente que visita la tienda, la gente que crea un carrito, el monto del carrito y la gente que no compra. Llevo un registro de clientes también: si una persona me compra más de una vez, puedo entrar y ver cuántas veces compraron y qué compraron. Además está vinculada con Correo Argentino y con Andreani; se maneja todo desde Tienda Nube, idealmente.

\noindent\textbf{Julián Fabbi:} Lo que me gustaría saber, y que solo lo tengo en la cabeza, es qué vende más. De todos los artículos que tengo cargados en la tienda, ¿cuál es el que más ventas tuvo? ¿Cuál es el que más interacciones tuvo? ¿Cuál es el más visto por el público?

\noindent\textit{--- Conocimiento del mercado e intuición ---}

\noindent\textbf{Tomás Gonçalves:} Es buena esa. ¿Cómo sabés hoy cuál se vende más?

\noindent\textbf{Julián Fabbi:} Yo sé que, por ejemplo, Jujutsu Kaisen vende un montón; es uno de los títulos que más vende. Blue Lock suele salir bastante también. Lo tengo en la cabeza: qué títulos salen, qué títulos van a salir. Por ejemplo, la temporada pasada estrenó un anime y las ventas de ese manga en particular subieron. Entonces yo sé que ahora está estrenando otro y las ventas del manga en el que está basado probablemente empiecen a subir. Pero lo tengo en la cabeza; es un dato que no puedo sacar de ningún lado, es conocer el mercado nada más.

\subsubsection*{Bloque 2 · Gestión de inventario y stock}

\noindent\textit{--- Indicadores de stock y alertas ---}

\noindent\textbf{Tomás Gonçalves:} ¿Tenés un indicador de cuándo tenés que reponer stock?

\noindent\textbf{Julián Fabbi:} Sí y no. En Tienda Nube tenés los productos con colores: rojo, amarillo y blanco. Lleva ese registro del stock que venís teniendo. No te avisa más nada que el colorcito, y cuando te quedás sin stock sí te llega un mensajito de ``productos sin stock''.

\noindent\textbf{Tomás Gonçalves:} O sea que es de repente: capaz que son las dos de la mañana y te enterás cuando te levantás al otro día.

\noindent\textbf{Julián Fabbi:} Sí. Y Mercado Libre es más o menos igual: te avisa que no tenés stock o que tenés poco, pero no te llega una notificación o una alerta. Por otro lado, es tan pequeño el negocio que tengo que no me puedo dar la comodidad de tener stock de sobra. Trabajo con muy poco stock de cada uno de los títulos.

\noindent\textit{--- Ciclo de reposición y dependencia del proveedor ---}

\noindent\textbf{Tomás Gonçalves:} Si te quedás sin stock de un producto, ¿cuánto tiempo pasa entre que te enterás y lo reponés?

\noindent\textbf{Julián Fabbi:} Depende del día. Yo recibo un listado todos los viernes y el sábado a las 14:00 es el último momento que tengo para hacer el pedido. Después de eso, la semana siguiente mi proveedor realiza las entregas. Si me entero justo el mismo viernes, es una semana. Si me llego a enterar el lunes, tengo que esperar al viernes y después una semana más.

\noindent\textbf{Julián Fabbi:} Dependo mucho del stock que tenga mi proveedor. Si un título está agotado, no hay forma de reponerlo hasta que la editorial lo reimprima. Y eso varía mucho: hay títulos que los reimprimen todos los meses y títulos que están agotados hace cinco meses y no salen.

\noindent\textit{--- Stock sobrante ---}

\noindent\textbf{Tomás Gonçalves:} ¿Te pasó que compraste stock y vendiste la mitad? ¿Qué hacés con lo que sobra?

\noindent\textbf{Julián Fabbi:} El stock que sobra lo guardo y de a poquito va saliendo. Hay algunos títulos que son más duros de sacar que otros, pero lo guardo.

\subsubsection*{Bloque 3 · Análisis de ventas y decisiones comerciales}

\noindent\textit{--- Publicidad y canales ---}

\noindent\textbf{Tomás Gonçalves:} ¿Hacés publicidad?

\noindent\textbf{Julián Fabbi:} Tengo publicidad en Mercado Libre y publicidad en Meta (Facebook e Instagram).

\noindent\textbf{Tomás Gonçalves:} ¿Tenés métricas de eso? Otra entrevistada nos decía que sabía que la gran mayoría de sus compras venían desde Instagram.

\noindent\textbf{Julián Fabbi:} No sé decirte de dónde vienen específicamente. Lo que sí: los que compran en Tienda Nube son los que vienen de Meta y TikTok, y los que compran en Mercado Libre son los que vienen de la publicidad en Mercado Libre, claramente. Pero los insights de la publicidad no los tengo todos juntos; tendría que buscar publicidad por publicidad a ver cómo le está yendo a cada una.

\noindent\textbf{Julián Fabbi:} Lo que sí te puedo decir es que por cinco mil pesos diarios tengo un promedio de doscientos clics en el link de Tienda Nube. Es una campaña que dura unas dos semanas y vienen tres compras. No son tantas, pero vienen. Donde más compras hay es en Mercado Libre: la publicidad es de tres mil pesos diarios, la tengo activa todo el mes, y promedio una compra por día, más o menos cada día y medio.

\noindent\textbf{Tomás Gonçalves:} ¿Recuperás la inversión?

\noindent\textbf{Julián Fabbi:} Sí. La publicidad la pago yo de mi bolsillo. Pero está dentro del plazo que yo esperaba: yo operaba con la idea de tener que mantener toda la publicidad entre uno y dos años hasta que esté andando solo. Todavía estamos en plan.

\noindent\textit{--- Estacionalidad y promociones ---}

\noindent\textbf{Tomás Gonçalves:} ¿Hay días que vendés más o es más esporádico?

\noindent\textbf{Julián Fabbi:} Es más esporádico. Hay algún día que tenés suerte y te caen dos compras de cien mil pesos, y otro día es una compra de veinte mil.

\noindent\textbf{Tomás Gonçalves:} ¿Coincide con Navidad o alguna fecha especial?

\noindent\textbf{Julián Fabbi:} Con Navidad sí, tuvimos muy buen período. Fue deliberadamente hecho así: para Navidad pusimos 15\% de descuento en toda la tienda, haciendo gancho con que la gente cobraba el aguinaldo. Las últimas dos semanas de diciembre fueron una explosión de compras. Y pasó lo mismo con la publicidad que metimos por la Feria del Libro: la primera semana de mayo pusimos el mismo descuento y también hubo una cantidad muy buena de compras.

\noindent\textbf{Julián Fabbi:} Ahora probamos con el aguinaldo de mitad de año y ese no salió tan bien. Dudo que haya sido por la publicidad, porque metí la misma o mil pesos menos, y no me parece proporcional a la diferencia de compras. Para mí fue porque mucha gente cobró el aguinaldo antes del cambio de mes, o porque estaban ahorrando para irse de vacaciones de invierno. No lo voy a saber si no le pregunto al público, pero son las dos ideas que tengo. Probablemente no vuelva a hacer la misma promoción a mitad de año.

\noindent\textit{--- Regulación de precios: Ley 25.542 ---}

\noindent\textbf{Tomás Gonçalves:} Tema precios: ¿los ajustás con frecuencia?

\noindent\textbf{Julián Fabbi:} No, no los puedo ajustar. En Argentina existe la Ley 25.542, de protección de la actividad librera. Uno de sus puntos es que los precios de los libros los define la editorial como precio de venta al público. Entonces, si me comprás a mí, pagás lo mismo que le pagarías a cualquier otra librería.

\noindent\textbf{Julián Fabbi:} La lógica es que todos vendamos libros al mismo precio. Está enfocada en que grandes empresas (supermercados, plataformas de e-commerce) no puedan hacer dumping vendiendo libros a pérdida, porque las librerías pequeñas no podrían competir con eso.

\noindent\textbf{Julián Fabbi:} Yo no puedo subir los precios porque no me van a comprar, y no los puedo bajar. Las promociones con descuento que hago están casi rozando la ilegalidad, pero no lo hago a precios de dumping: sigo teniendo una ganancia, es un poco menor, y lo hago por un período corto de tiempo como estrategia publicitaria para que me conozcan. Los precios los definen las editoriales y de vez en cuando actualizan: este año Ivrea actualizó una vez, Planeta y Penguin actualizaron dos veces.

\subsubsection*{Bloque 4 · Puntos de dolor y visibilidad del negocio}

\noindent\textbf{Tomás Gonçalves:} ¿Sentís que tenés visibilidad clara de lo que está pasando día a día en Tienda Nube o en Mercado Libre?

\noindent\textbf{Julián Fabbi:} Tengo las métricas de que la gente entró a la tienda, de que tomó este producto al carrito. Pero me gustaría ver si el buscador de la tienda registra qué buscan los usuarios. Qué títulos entran a ver de los que hay disponibles, qué títulos buscan. Eso estaría bueno tenerlo.

\subsubsection*{Bloque 5 · Validación de la propuesta de valor}

\noindent\textit{--- Confianza y disposición a pagar ---}

\noindent\textbf{Tomás Gonçalves:} Si tuviéras una herramienta que todos los días te dice cosas como ``conviene reponer X producto antes del fin de semana'', ``lanzá una promo mañana porque es una fecha clave'' o ``pausá los ads del producto porque casi no queda stock'', ¿cuánto valor le darías? ¿Cuánto pagarías como plugin para Tienda Nube?

\noindent\textbf{Julián Fabbi:} Confiaría en un plugin así, sí. Actualmente estoy pagando creo que veintisiete mil pesos de Tienda Nube del plan más barato, más un 6\% de retenciones por venta. Hay planes más caros de hasta doscientos cincuenta mil pesos, que yo dije: ¿qué? Lo que estoy usando me sirve, me alcanza y me sobra.

\noindent\textbf{Julián Fabbi:} Si lo tuviera que sumar al costo que vengo pagando de Tienda Nube, yo sumaría entre cinco mil y diez mil pesos. No mucho más. Pero es relativo al plan que pago. Si estuviese pagando un plan más alto de ciento cincuenta o doscientos mil pesos, capaz te digo: sumále treinta mil al plugin. Lo jugaría mucho en base a lo que pago en el momento.

\noindent\textit{--- Funcionalidad más deseada ---}

\noindent\textbf{Tomás Gonçalves:} ¿Qué es lo que más necesitás ahora que te ayudaría un montón?

\noindent\textbf{Julián Fabbi:} Lo que más me gustaría saber es cuáles son los títulos que más entran a ver y cuáles son los que menos entran a ver. Porque de los que más entran a ver, digo: traigo más de esto. De los que menos entran a ver: hago contenido en redes sociales para esos títulos justamente. Sería lo que más me interesaría saber en este momento.

\subsubsection*{Bloque 6 · Cierre}

\noindent\textbf{Tomás Gonçalves:} ¿Hay algo que no te pregunté que te gustaría decir?

\noindent\textbf{Julián Fabbi:} No, creo que está todo.

\noindent\textbf{Julián Fabbi:} Vos avisame si necesitás más. Si necesitamos volver a charlar o información propia de la tienda, no hay problema.

\noindent\textbf{Tomás Gonçalves:} Mil gracias, Julián. Cuando lancemos la aplicación, te damos el acceso al tier más alto.

\begin{center}\textit{--- Fin de la entrevista ---}\end{center}
